\documentclass[12pt]{article}

\usepackage{geometry}
\usepackage{titling}
\usepackage{amsmath,amssymb}
\usepackage{fontspec}
\usepackage{unicode-math}
\usepackage{array}
\usepackage{colortbl}
\usepackage[dvipsnames]{xcolor}
\usepackage[most]{tcolorbox}

\geometry{a4paper, top=2.5cm, bottom=2.5cm, left=2.6cm, right=2.5cm}
\setlength{\droptitle}{-3cm}

\setmonofont{DejaVu Sans Mono}
\setmainfont{Libertinus Serif}
\setsansfont{Libertinus Sans}
\setmathfont{STIXTwoMath-Regular.otf}

\newcommand{\MVec}[1]{\mathbf{#1}}

\title{What Standard Deviation Actually Means}
\author{Devi Prasad M}
\date{September 2026}

\begin{document}
\pagecolor{yellow!30!brown!40}

\maketitle

\section*{Introduction}

Standard deviation is often introduced as ``the typical amount by
which entries differ from their average.'' That phrasing is
intuitive but can be misleading --- it makes one to compare
$\mathbf{std}(\MVec{x})$ against something like the
average deviation size. These two things don't
generally agree. Therefore, it is more accurate to describe
it in geometric terms, rather than using a naive language.

\begin{quote}
\textbf{std($\MVec{x}$) is the unique uniform magnitude that,
spread evenly across every entry, reproduces the same total
squared deviation as the actual (uneven) data.}
\end{quote}

It is defined from squared quantities, and is disproportionately
sensitive to entries that deviate the most.

\section*{An Example}

Consider the vector
\[
x = (-0.46,\ -0.51,\ -0.29,\ 0.18,\ 1.22,\ -0.79,\ -1.27,\ 1.92),
\]
which has $\mathbf{avg}(\MVec{x}) = 0$ and
\begin{align*}
    \mathbf{std}(\MVec{x}) &= \sqrt{\dfrac{\sum_i (x_i - \mathbf{avg}(\MVec{x}))^2}{n}}\\[2pt]
                           &= \sqrt{\frac{(-0.46)^2 + (-0.51)^2 + \dots + (-1.27)^2 + (1.92)^2}{8}}\\[2pt]
                           &= \sqrt{\dfrac{8}{8}}\\
                           & = 1
\end{align*}
exactly.

\noindent The simple exercise of averaging the absolute
deviations, gives Mean Absolute Deviation (MAD):
\[
    \mathbf{MAD} = \dfrac{0.46+0.51+0.29+0.18+1.22+0.79+1.27+1.92}{8} = 0.83,
\]
not $1$. The gap exists because one entry, $1.92$,
contributes $1.92^2 = 3.69$ to the sum of squares ---
almost half the total --- while the other seven together
contribute the rest. Squaring inflates the influence of large
deviations, so $\mathbf{std}$ ends up pulled above the ``typical'' value
one would eyeball from the data.

\section*{What $\mathbf{std}$ Actually Means}

if $\overline{\MVec{x}} = \MVec{x} - \mathbf{avg}(\MVec{x})\mathbf{1}$
is the de-meaned vector, then
\[
\mathbf{std}(\MVec{x}) = \dfrac{\|\overline{\MVec{x}}\|}{\sqrt n}
\]
is the value $r$ such that a vector with every entry equal
to $r$ in magnitude has the same Euclidean length as $\overline{\MVec{x}}$.
It does \textbf{not} claim that individual entries
resemble $r$ --- only that total squared deviation is conserved.

\subsection*{Finance analogy}
A \textbf{fixed deposit} has $\mathbf{std} = 0$ (no fluctuation, ever).
A \textbf{real stock} fluctuates unevenly --- big swings some periods,
small ones others. The idea of $\mathbf{std}$ answers the question:

\begin{quote}
\emph{If the same total risk (or Euclidean length)
were spread perfectly evenly instead of lumped unevenly,
what size would each period's swing have to be?}
\end{quote}

If $\mathbf{std} = 10$ (in percentage points) around a mean
return of $X\%$, the equivalent perfectly regular ``seesaw''
alternates between $X+10$ and $X-10$ every period --- not because
the real stock actually moves that predictably, but because that
regular seesaw is the flat instrument carrying identical total risk.

\section*{Conclusion}
It is important to understand that $\mathbf{std}$ is not
``how much things usually differ'' --- it is ``the flat,
constant-swing version of your data that has the same total squared
spread.''

When deviations are roughly equal in size (like the seesaw),
that flat version matches intuition well. When one or two entries
dominate (like $1.92$ in the example above), $\mathbf{std}$
still measures the real total spread correctly --- it
no longer looks like an ``average-sized'' deviation,
because it isn't one.

When we say a dataset has a standard deviation of $r$,
\begin{itemize}
    \item it does not mean most entries are equal to $r$.
    \item It does not mean $r$ is the simple average of the distances.
    \item It means that if all fluctuations were smoothed out
          into an equal-sized, regular heartbeat of $\pm r$,
          the total combined impact (variance) would remain identical.
\end{itemize}

When we say $\mathbf{std}(x) = 1.0$ is ``typical,"
we actually mean:
\begin{quote}
    If we replaced all these varying deviations with a single
    uniform magnitude that preserves the total Euclidean length
    of the vector, that magnitude would be $1.0$.
\end{quote}

\section*{Additional Notes}
\noindent If $\mathbf{std}(\MVec{x})$ is the constant $r \ge 0$
such that a vector $v \in \mathbb{R}^n$ with all entries
of magnitude $r$ (i.e.\ $|v_i| = r$ for every $i$, signs arbitrary),
it has the same norm as the de-meaned vector $\overline{\MVec{x}}$.

Check it: $\|v\|^2 = \sum_i v_i^2 = \sum_i r^2 = n r^2$,
so $\|v\| = r\sqrt{n}$. Setting $\|v\| = \|\overline{\MVec{x}}\|$:
\[
    r\sqrt{n} = \|\overline{\MVec{x}}\|
        \quad\Longrightarrow\quad r
        = \frac{\|\overline{\MVec{x}}\|}{\sqrt{n}}
        = \mathbf{std}(\MVec{x}).
\]

$\mathbf{std}(\MVec{x})$ is the unique uniform magnitude
that reproduces $\|\overline{\MVec{x}}\|$ across $n$ coordinates.

\end{document}